\begin{figure}[!t]
    \centering
    \begin{minipage}{0.47\linewidth}
        \centering
        \includegraphics[width=\linewidth]{figures/bench/wan-bench.pdf}
        \caption{The dimensions covered in \method-Bench.}
        \label{vanbench-struct}
    \end{minipage}\hfill
    \begin{minipage}{0.47\linewidth}
        \centering
        \includegraphics[width=\linewidth]{figures/bench/prompt_dist.pdf}
        \caption{Distribution of the number of prompts across various dimensions of \method-Bench.}
        \label{prompt_dist}
    \end{minipage}
\end{figure}



\subsection{Benchmarks}

Existing video generation evaluation methods, such as Fréchet Video Distance (FVD)~\citep{2018Towards} and Fréchet Inception Distance (FID)~\citep{heusel2017gans}, lack alignment with human perception. We propose an automated, comprehensive, and human-aligned \method-Bench for evaluating video generation models. 


\method-Bench is composed of three core dimensions: dynamic quality, image quality, and instruction following with 14 fine-grained metrics. We design specific scoring algorithm for each dimension, using traditional detectors for the simple tasks (\eg, object detection) and multi-modal Large Language Models (MLLMs) for the complex tasks. 


\textbf{Dynamic quality}. Dynamic quality aims to evaluate the quality and stability of the model in non-static scenarios.

\textit{Large motion generation.} To measure the upper limit of motion generation, prompts are designed to encourage significant motions. 
%
We calculate the optical flow of the generated video using RAFT~\citep{2021RAFT} and assess the motion scores by the magnitude of the normalized optical flow.


\textit{Human artifacts.}
Existing evaluation methods lack the ability to detect AI-generated artifacts.
%
Thus, we train a YOLOv3-based model~\cite{2018YOLOv3} on 20,000 AI-generated images annotated by humans to identify the locations of artifacts.
%
The artifact scores take into account the predicted likelihood, the bounding boxes, and the corresponding duration.


\textit{Physical plausibility \& smoothness}. To evaluate physical plausibility, we design prompts to generate physics-related movements (\eg, ball bouncing, object interactions, water flow) and employ Qwen2-VL~\citep{qwenvl}, leveraging its extensive real-world physics knowledge. 
%
Through a video question-answer model, we detect violations of physical laws, such as object clipping, unrealistic collisions, or gravity-defying movements.
To evaluate smoothness, prompts are crafted to include complex motions, with Qwen2-VL identifying artifacts to assess motion fluidity.



\textit{Pixel-level stability}.
Pixel-level stability is used to determine if noise frequently appears in the video. We employ optical flow to identify static regions and calculate frame-wise differences within those areas. 
For stable videos, these differences within static regions are relatively minimal.


\textit{ID consistency.} Three kinds of sub-dimensions are considered: human consistency, animal consistency, and object consistency covering different collections of prompts. Frame-level DINO \citep{2021Emerging} features are extracted to compare the frame-wise similarity as the consistency score.

\textbf{Image quality.} The purpose of the image quality dimension is to evaluate the visual quality and aesthetic perception.


\textit{Comprehensive image quality.} Measuring the visual quality requires an evaluation of the fidelity and aesthetics. 
For fidelity, we employ MANIQA \citep{2022MANIQA}, which effectively detects blur and artifacts, ensuring precise clarity assessment. 
For aesthetics, we utilize two established models: the LAION-based aesthetic predictor~\citep{LAION-AI_aesthetic-predictor_2022} and MUSIQ~\citep{2021MUSIQ}. We compute the final image quality score by averaging these three evaluators.


\textit{Scene generation quality.}
We evaluate two key aspects: frame-wise scene consistency and scene-text alignment. For consistency, we measure CLIP similarity \citep{radford2021learning} across consecutive frames to ensure temporal coherence. 
For alignment, we compute CLIP similarity between frames and their corresponding text to verify semantic accuracy. 
We compute the quality score by weighted averaging these metrics.


\textit{Stylization.} We employ Qwen2-VL to evaluate the artistic video generation capability through frame-level question answering.


\textbf{Instruction following.} Instruction following assesses the model's ability to comprehend and execute textual instructions.


\textit{Single object \& multiple objects \& spatial positions.} 
We employ Qwen2-VL to predict object categories, counts, and spatial relationships within videos. 
The evaluation score represents the average number of frames that accurately correspond to the given prompts.



\textit{Camera control.}
We evaluate five camera movement styles: panning, lifting, zooming in/out, aerial shots, and tracking shots by 130 tailored prompts. For panning, lifting, and zooming, we analyze optical flow via RAFT to assess movement style. 
For complex aerial and tracking shots, we employ Qwen2-VL to infer camera movement through video question-answering.

\textit{Action instruction following.} We categorize motions into human (\eg, running), animal (\eg, crawling), and object movements (\eg, flying). 
As evaluating these motions requires prior knowledge of their execution, we employ Qwen2-VL by providing it with keyframes. 
We query the model for alignment of the action, completion of the action, and presence of artifacts to comprehensively evaluate alignment between text and video.

\textbf{Human feedback guided weighting strategy.}\label{hf-align} We utilize human feedback to assign a comprehensive video quality score, considering that each dimension has a different impact on user preferences rather than simply averaging them.
We collect over 5,000 pairwise comparisons of videos generated by different models including \method. Users evaluate pairs based on a given text, indicating their preference and assigning approximate scores. 
The dimension preference weighting is derived from the Pearson correlation between model-generated scores and human-rated scores, serving as a weighting factor in the final score calculation. 